\section{Introdução}

% Nos últimos anos, a comercialização de alimentos destinados à alimentação escolar tem ganhado relevância crescente no Brasil, especialmente no contexto do Programa Nacional de Alimentação Escolar (PNAE). Em 2025, a Lei nº 15.226 ampliou de 30\% para 45\% o percentual mínimo dos recursos do programa que devem ser destinados à aquisição direta de gêneros alimentícios da agricultura familiar, com vigência a partir de 2026~\cite{fnde_pnae_2025}. Além disso, em fevereiro de 2026, o governo federal anunciou um reajuste de 14,35\% nos repasses da merenda escolar, elevando o investimento total no programa para R\$ 6,7 bilhões~\cite{agencia_brasil_merenda_2026}. Tais medidas ampliam significativamente a importância econômica desse mercado e intensificam a participação de cooperativas, associações e pequenos produtores como agentes fornecedores.

% Nesse cenário, a questão central não se resume apenas ao processo de compra por parte do poder público, mas também às decisões estratégicas dos próprios vendedores. Cooperativas e fornecedores da agricultura familiar frequentemente ofertam múltiplos produtos (como arroz, feijão, hortaliças, frutas e derivados) para diferentes redes de ensino ou chamadas públicas, cujas condições de aquisição podem variar conforme restrições orçamentárias, custos logísticos, sazonalidade e preços de referência. Em particular, compradores distintos podem estar dispostos a aceitar valores diferentes para um mesmo item, o que torna a definição de preços uma decisão não trivial por parte do vendedor.

% Surge, então, um dilema econômico natural: \blue{preços mais elevados aumentam a receita unitária, mas podem reduzir o conjunto de compradores dispostos a adquirir determinado produto; por outro lado, preços mais baixos ampliam o número de vendas possíveis, porém diminuem o ganho obtido em cada transação.} Assim, do ponto de vista do vendedor, o problema consiste em definir preços para um conjunto de itens de forma a equilibrar essas forças antagônicas e maximizar a receita total obtida.

% Esse tipo de situação se insere em uma classe mais ampla de problemas conhecidos como \textit{problemas de precificação}, nos quais um vendedor deve determinar os preços de seus produtos de modo a maximizar a receita total obtida com as vendas. Em contextos práticos, essa decisão é particularmente desafiadora, pois a definição do preço de cada item afeta diretamente tanto o valor arrecadado por unidade vendida quanto o conjunto de compradores dispostos a adquiri-lo. Dessa forma, a precificação adequada exige um balanceamento entre competitividade e rentabilidade, tornando-se um problema de otimização de grande relevância teórica e aplicada.

% De maneira geral, problemas de precificação envolvem determinar os preços dos produtos de um vendedor de forma que sua receita (entendida como a soma dos valores obtidos com todas as vendas realizadas) seja a maior possível. Em termos intuitivos, \blue{preços mais altos podem aumentar o retorno por unidade comercializada, mas tendem a restringir o conjunto de compradores interessados; por outro lado, preços mais baixos podem ampliar o volume de vendas, embora reduzam o ganho em cada transação.} Esse \textit{trade-off} aparece em diversos setores econômicos e tem sido amplamente estudado sob diferentes perspectivas, dentre as quais se destacam as formulações matemáticas e o desenvolvimento de algoritmos exatos.\red{Acho que cabe algumas referências aqui}


% \blue{Daqui pra baixo eu devo falar sobre algumas trabalhos na área de pricing. Talvez iniciar falando sobre algumas variantes do problema, o que já se tem para eles (exatos, heurísticas) e levar ao artigo que define o problema formalmente. Depois disso eu devo de alguma forma falar sobre o trabalho do Sandro, falando sobre as formulações dele e o que eu pretendo fazer aqui neste trabalho}


% MARK

Relações comerciais estão presentes no cotidiano de todos. De modo geral, um vendedor expõe seus itens por um determinado valor, e um comprador, ao estar disposto a desembolsar tal quantia, pode adquiri-los. Em um mercado cada vez mais competitivo, surge uma questão central para o vendedor: como precificar seus produtos da melhor maneira possível? Nesse contexto, o problema não se limita apenas à definição de um preço, mas envolve também a compreensão do comportamento dos consumidores em diferentes cenários. Por exemplo, um comprador pode optar por adquirir um produto alternativo ao encontrar outro que ofereça a mesma funcionalidade por um preço mais acessível.

Surge, então, um dilema econômico: preços mais elevados aumentam a receita unitária, mas podem reduzir o conjunto de compradores dispostos a adquirir determinado produto; por outro lado, preços mais baixos ampliam o número de vendas possíveis, porém diminuem o ganho obtido em cada transação. Assim, do ponto de vista do vendedor, o problema consiste em definir preços para um conjunto de itens de forma a equilibrar essas forças antagônicas e maximizar a receita total obtida.

Esse tipo de situação se insere em uma classe mais ampla de problemas conhecidos como \textit{Problemas de Precificação} (PP), nos quais um vendedor deve determinar os preços de seus produtos de modo a maximizar a receita total obtida com as vendas. Essa decisão é desafiadora, pois o preço de cada item influencia simultaneamente o valor arrecadado por unidade e o conjunto de consumidores dispostos a adquiri-lo, evidenciando a necessidade de um equilibrio entre competitividade e rentabilidade~\cite{smith1986}.

Problemas de precificação estão presentes em diversos setores da economia e têm sido amplamente investigados na literatura, devido à sua elevada relevância prática e aplicabilidade em contextos reais~\cite{aggarwal2004,venkatesan2005,shmuel1984,shmuel1987,smith1986}. Nesse cenário, Rusmevichientong~\cite{rusmevichientong2006} formalizou o conceito de problemas de precificação não paramétrica com múltiplos itens e compradores, propondo diferentes variantes para modelar distintos comportamentos de compra. Dentre essas variantes, destacam-se o \textit{Problema da Compra Máxima} (PCmax), no qual o comprador seleciona o produto viável de maior preço, e o \textit{Problema da Compra por Preferência} (PCrank), em que o comprador escolhe, dentre os itens viáveis, aquele que ocupa a posição mais alta em sua ordem de preferência.

% Entre as diferentes variantes propostas por Rusmevichientong, destaca-se o \textit{Problema da Compra Mínima} (PCmin), definido como segue. Seja um conjunto de itens $I$, um conjunto de compradores $B$ e uma matriz de valorações $V$, em que cada elemento $v_{ib} \in V$ representa a valoração do comprador $b \in B$ para o item $i \in I$. O objetivo do problema é determinar um conjunto de preços $P$, no qual cada elemento $p_i$ corresponde ao preço atribuído ao item $i$, de forma a maximizar a receita total do vendedor. Além disso, assume-se que o comportamento dos compradores é conhecido: cada comprador adquire um único item viável, escolhendo sempre aquele de menor preço. Caso nenhum item seja viável, nenhuma compra é realizada.

% Entre as diferentes variantes propostas por Rusmevichientong, destaca-se o \textit{Problema da Compra Mínima} (PCmin), definido como segue. Seja um conjunto de itens $I$, um conjunto de compradores $B$ e uma matriz de valorações $V$, em que cada elemento $v_{ib} \in V$ representa a valoração do comprador $b \in B$ para o item $i \in I$. O objetivo do problema é determinar um conjunto de preços $P$, no qual cada elemento $p_i \in P$ corresponde ao preço atribuído ao item $i$, de forma a maximizar a receita total do vendedor. Assume-se, adicionalmente, que o comportamento dos compradores é conhecido e segue uma regra determinística: cada comprador $b \in B$ adquire, dentre os itens de seu interesse, o item de menor preço cujo valor não exceda sua valoração. Isto é, o comprador escolhe um item $i \in I$ tal que $p_i \leq v_{ib}$ e para qualquer $j \in I$, $p_i \leq p_j$. Caso não exista item viável, o comprador não realiza compra.

Entre as diferentes variantes propostas por Rusmevichientong, destaca-se o \textit{Problema da Compra Mínima} (PCmin), definido como segue. Seja um conjunto de itens $I$, um conjunto de compradores $B$ e uma matriz de valorações $V$, em que cada elemento $v_{ib} \in V$ representa a valoração do comprador $b \in B$ para o item $i \in I$. O objetivo do problema é determinar um conjunto de preços $P$, no qual cada elemento $p_i \in P$ corresponde ao preço atribuído ao item $i$, de forma a maximizar a receita total do vendedor. Assume-se, adicionalmente, que o comportamento dos compradores é conhecido e segue uma regra determinística: cada comprador $b \in B$ adquire, dentre os itens de seu interesse, o item de menor preço cujo valor não exceda sua valoração. Isto é, o comprador escolhe um item $i \in I$ tal que $p_i \leq v_{ib}$ e, para qualquer item $j \in I$ que satisfaça $p_j \leq v_{jb}$, tem-se $p_i \leq p_j$. Caso não exista item viável, o comprador não realiza compra.

Essa hipótese de comportamento torna o PCmin particularmente relevante do ponto de vista prático, uma vez que, em diversos contextos reais, consumidores tendem a optar pelo item de menor preço dentre aqueles que consideram aceitáveis ou de seu interesse. Embora fatores como marca, qualidade percebida e preferências individuais também possam influenciar a decisão de compra, o critério de menor preço frequentemente exerce papel determinante, especialmente em mercados com produtos substitutos ou de baixa diferenciação.

\begin{figure}[H]
\centering

\begin{subfigure}[t]{0.45\textwidth}
    \centering
    \conjuntoItensValorados
    \caption{Relações de valoração entre itens e compradores, representadas pelas arestas ponderadas da matriz $V$.}
\end{subfigure}
\hfill
\begin{subfigure}[t]{0.45\textwidth}
    \centering
    \conjuntoItensPrecificados
    \caption{Preços atribuídos aos itens e seleção induzida dos compradores segundo a regra do PCmin.}
\end{subfigure}

\caption{Exemplo ilustrativo de uma instância do PCmin, com as valorações entre itens e compradores e a seleção induzida pelos preços atribuídos.}

\label{fig:exemploPCmin}
\end{figure}

A Figura~\ref{fig:exemploPCmin} apresenta um exemplo ilustrativo, não necessariamente ótimo, de uma instância do PCmin. Na Figura~\ref{fig:exemploPCmin}a, são representadas as relações de valoração entre os itens e os compradores por meio de um grafo bipartido, no qual cada aresta ponderada indica a valoração $v_{ib}$ atribuída pelo comprador $b$ ao item $i$. Já na Figura~\ref{fig:exemploPCmin}b, são mostrados os preços atribuídos aos itens e as escolhas induzidas pela regra do problema. Nesse caso, cada comprador adquire o item de menor preço dentre aqueles de seu interesse cujo valor não excede sua respectiva valoração. Assim, observa-se que o comprador $b_1$ seleciona o item $i_1$, o comprador $b_2$ seleciona o item $i_4$, o comprador $b_3$ seleciona o item $i_3$, enquanto o comprador $b_4$ não realiza compra, uma vez que não há item viável disponível para ele.


Dentre os trabalhos mais recentes relacionados ao PCmin, destaca-se o de Catabriga \cite{catabriga2025}, que propõe diferentes formulações de Programação Linear Inteira (PLI) e Inteira Mista (PLIM) para o problema, além de explorar técnicas de fortalecimento dos modelos e avaliar seu desempenho computacional em diversas classes de instâncias. Seus resultados evidenciam o potencial dessas abordagens para a resolução do problema, bem como os desafios ainda existentes em termos de escalabilidade e eficiência.

Neste contexto, o presente trabalho tem como objetivo avançar nessa linha de pesquisa, investigando possíveis melhorias nas modelagens e implementações propostas. Em particular, buscamos desenvolver formulações mais eficientes e explorar estratégias que possam reduzir o tempo de execução. Adicionalmente, pretendemos implementar e avaliar abordagens heurísticas voltadas à resolução de instâncias de grande porte, nas quais métodos exatos podem se tornar computacionalmente inviáveis, contribuindo assim para o aprimoramento das abordagens existentes na literatura.
